% LTeX: language=en-GB
\chapter*{Summary}
\addcontentsline{toc}{chapter}{Summary}

In Turin (north-west Italy), a research line between \textit{Politecnico di Torino}, the \textit{A.O.U. Città della Salute e della Scienza di Torino} hospital complex and \textit{FEV Italia S.r.l.} has studied since 2020 how drones could support the ambulance-based transport of medical supplies between hospitals. If a drone fails in flight, however, its payload may detach and fall onto the city, where the impact could, among other consequences, destroy urgently needed medical goods and disperse potentially hazardous substances in the environment. No device currently exists to protect such a payload, nor a method to design one for a given mission. An airbag could fulfil this role, but sizing it requires two quantities: the impact velocity of the payload and the largest acceleration it can withstand without damage.

This thesis addresses the problem on two levels. An outer design loop formalises a safety-first principle, in which the payload and the drone are treated as separate entities and the protection of the payload is designed at the start rather than added at the end, so that it supplies the drone design with inputs such as airbag mass, volume, inflator type and inflation time. Within this loop, an inner design loop forms the workflow delivered by the thesis: it establishes the two quantities above and returns the lightest airbag that keeps the payload below its allowable acceleration. Its first five steps, from mission analysis to design optimisation, were demonstrated on a reference delivery mission between two hospitals in Turin, flown largely above the river Po and carrying 5\,ml borosilicate glass phials, while the remaining three steps depend on the geometry of the container and the drone and are left to future work.

The first step of the workflow, a mission analysis, established the impact conditions: a Monte Carlo simulation of the fall along the whole route gave, at its 98.5th percentile, a design impact of 7.5\,m/s horizontally and 25\,m/s vertically. The next two steps characterised the payload: compression tests on 76 phials from two producers showed that nominally identical phials behaved differently, so only the batch with the smaller scatter was retained, and its 5th-percentile failure load of 98.7\,N became the target for a finite element model of the same test. With a glass tensile strength of 60\,MPa, the model matched this load within 2\% and reproduced the two most common failure modes of that batch. In the fourth step, the calibrated phial was placed in a simulated sled test, in which an assumed container, lined with a foam sponge and a plastic tray, was subjected to an acceleration pulse emulating the impact at the resultant speed of 26.1\,m/s. Loaded perpendicular to its axis, the most critical of the directions tested, the phial withstood a peak of 225\,$g$ and cracked at 230\,$g$, so 225\,$g$ was taken as the allowable acceleration. The final step brought both strands together: an analytical airbag impact model, extended to account for horizontal velocity, fabric porosity and bursting, was embedded in a differential evolution optimiser. Given a 1\,kg payload, both velocity components and an acceleration limit of 187.5\,$g$ (the allowable divided by a safety factor of 1.2), the optimiser selected the same design for every landing-speed limit between 0.5 and 2.0\,m/s: a single-ply cylinder of silicone-coated Nylon 6,6 inflated with helium, 0.522 to 0.774\,m in diameter depending on that limit, with fabric and gas together below 0.4\,kg.

The workflow thus sizes the inflated airbag for a specific payload and impact scenario. Its outputs point to a stored-gas helium inflator, since helium cannot be produced pyrotechnically, although the feasibility and mass of the pressurised gas containers, which the optimisation does not include, remain to be verified. The design is only as reliable as its inputs: a single phial geometry and batch were modelled, every optimum lies on at least one bound of the search space, and the airbag model has not yet been checked against physical tests. Once the capsule geometry is fixed, the optimised airbag should be validated through simulations of the real assembly and drop tests, and the workflow could be extended to a reliability-based design, combining the scatter in glass strength with that of the impact velocities, or to an airbag working in tandem with a parachute.